% ---------------------------------------------------------------------
% Section 4: Relationships among information orders.
% Drafted by Claude from Kim (2023, JET), per the authors' instruction,
% with \Claude{} comments flagging items for the authors to verify.
% Kim's definitions and results are restated in this paper's notation:
% states theta, signal processes X (better) and Y (worse); Kim's F and
% G correspond to our L_X and L_Y. Static comparison: the action/period
% arguments are suppressed throughout this section.
% ---------------------------------------------------------------------
\section{Relationships Among Information Orders}\label{sec:orders}

In this section, we relate the LR-better order to three information
orders from the literature: Blackwell's garbling order
\citep{blackwell1951comparison}, Lehmann's accuracy order
\citep{lehmann1988comparing}, and the monotone quasi-garbling order of
\citet{kim2023comparing}. Throughout this section, we compare the
signal processes in a single period and suppress the action argument;
in the dynamic model, the comparisons are assumed to hold for each
action. We write $F_\iX(x|\theta)$ for the cumulative distribution
function associated with the likelihood $L_\iX(x|\theta)$, and
similarly for $\iY$.

We begin by recalling the three orders. \citet{kim2023comparing}
introduces the monotone quasi-garbling order by relaxing Blackwell's
requirement that the added noise be independent of the state, while
restricting the state dependence to be \emph{reversely monotone} in
the sense of FOSD.

\begin{definition}[Garbling and monotone quasi-garbling]\label{def:BW-MQG}\quad
\begin{enumerate}
\item \citep{blackwell1951comparison} $\iX$ is \emph{Blackwell-better}
  than $\iY$ ($\iX \BW \iY$) if there exists a conditional probability
  distribution $P(y|x)$ such that
  $L_\iY(y|\theta) = \int P(y|x) L_\iX(x|\theta) \dif x$ for all
  $y \in Y$ and $\theta$.
\item \citep[Definition~3]{kim2023comparing} $\iX$ is a \emph{monotone
  quasi-garbling improvement} over $\iY$ ($\iX \MQG \iY$) if there
  exists a conditional probability distribution $P(y|x,\theta)$ such
  that (a) $L_\iY(y|\theta) = \int P(y|x,\theta) L_\iX(x|\theta)
  \dif x$ for all $y \in Y$ and $\theta$, and (b)
  $P(\,\cdot\,|x,\theta_1) \FOSD P(\,\cdot\,|x,\theta_2)$ for all
  $x \in X$ and $\theta_2 \geq \theta_1$.
\end{enumerate}
\end{definition}

\begin{definition}[Lehmann accuracy]\label{def:Lehmann}
\citep{lehmann1988comparing, kim2023comparing} $\iX$ is \emph{more
Lehmann accurate} than $\iY$ ($\iX \Leh \iY$) if, for all $y \in Y$,
the function
$\phi(\theta; y) = \sup\{x \,|\, F_\iX(x|\theta) \leq
F_\iY(y|\theta)\}$ is increasing in $\theta$.
\end{definition}

Lehmann's order has a simple interpretation: $\phi(\,\cdot\,; y)$ maps
a signal $y$ from $\iY$ to the signal from $\iX$ at the same
conditional quantile; $\iX$ is more accurate if this
quantile-matching map moves in the same direction as the state.
\citet{lehmann1988comparing} shows that, when payoffs and likelihoods
have appropriate monotone structure, more accurate experiments are
more valuable in single-decision problems; \citet{persico2000information}
and \citet{quah2009comparative} extend the reach of this result to
broader classes of monotone decision problems.
\citet{kim2023comparing} shows that, under the MLR property, Lehmann's
order is \emph{equivalent} to the monotone quasi-garbling order, which
gives Lehmann's condition a garbling interpretation.

The definitions of the four orders differ only in what they require of
the garbling kernel $P$ (or, for Lehmann's order, in avoiding a kernel
altogether): Blackwell requires state independence, the LR-better
order requires the state dependence to be reversely monotone in the LR
sense, and monotone quasi-garbling requires it to be reversely
monotone in the weaker FOSD sense. The following proposition collects
the resulting relationships.

\begin{proposition}[Relationships among the orders]\label{prop:orders}\quad
\begin{enumerate}
\item $\iX \BW \iY$ implies $\iX \LRI \iY$, and $\iX \LRI \iY$ implies
  $\iX \MQG \iY$. Neither implication can be reversed.
\item If $L_\iX$ satisfies the MLR property, then $\iX \MQG \iY$ if
  and only if $\iX \Leh \iY$ \citep[Theorem~1]{kim2023comparing}.
\item Without the MLR property on $L_\iX$, neither $\iX \BW \iY$ nor
  $\iX \MQG \iY$ implies $\iX \Leh \iY$
  \citep[Proposition~4]{kim2023comparing}.
\end{enumerate}
\end{proposition}

\begin{proof}{Proof of (i).}
If $\iX \BW \iY$ with a state-independent kernel $P(y|x)$, then
$P(y|x,\theta) = P(y|x)$ satisfies condition (ii) of
Definition~\ref{def:LRI} with equality, so $\iX \LRI \iY$. If
$\iX \LRI \iY$, then condition (ii) of Definition~\ref{def:LRI} states
that $P(\,\cdot\,|x,\theta_1) \LR P(\,\cdot\,|x,\theta_2)$ for
$\theta_2 \geq \theta_1$; since LR dominance implies FOSD, the same
kernel satisfies Definition~\ref{def:BW-MQG}(ii), so $\iX \MQG \iY$.
For the strictness of the first implication, the Normal example of
Section~\ref{sec:examples} with a strictly decreasing bias function
$d$ provides an LR-better garbling that is not a Blackwell garbling.
\Claude{Strictness claims to be completed: (a) for
$\BW \Rightarrow \LRI$, we should verify formally that the Normal
example with decreasing $d$ is not Blackwell-comparable (intuitively
the noise is state dependent, but a different state-independent kernel
could conceivably generate the same $L_\iY$---this needs an argument);
(b) for $\LRI \Rightarrow \MQG$, we need an example of a garbling
kernel that is reversely FOSD-monotone but not reversely
LR-monotone---and, ideally, such that no LR-monotone kernel generates
the same $\iY$. Your note at the end of the Uniform example in the
source file (mixtures of deterministic garblings, and uniform
garblings on $[x, x+d(\theta)]$) suggests candidate constructions.
Please advise which example you would like to develop.}
\Halmos
\end{proof}

Under the MLR property, then, the orders are nested:
\begin{equation}
\iX \BW \iY \ \Rightarrow\ \iX \LRI \iY \ \Rightarrow\
\iX \MQG \iY \ \Leftrightarrow\ \iX \Leh \iY \ .
\label{eq:nesting}
\end{equation}
\citet{kim2023comparing} shows that, in \emph{static} monotone
decision problems, the monotone quasi-garbling order characterizes the
value of information: it is sufficient for one signal process to yield
a higher expected payoff for all decision makers with monotone
decision problems (his Theorem~2), and it is also necessary when the
class of decision problems is rich enough to include all simple
hypothesis-testing problems (his Theorem~3 and Corollary~1). In view
of \eqref{eq:nesting}, one might hope that the same order would
suffice for \emph{dynamic} monotone decision problems. The reason we
work with the stronger LR-better order is the same reason the LR order
replaces FOSD throughout Section~\ref{sec:model}: an FOSD comparison
between garbling distributions need not survive Bayesian updating and
state transitions, while the LR comparison does. The two-state
example of Section~\ref{sec:tech-example} illustrates what can go
wrong when the reversely monotone structure fails; whether the
monotone quasi-garbling order itself is sufficient for value
comparisons in monotone POMDPs---or whether the gap between
$\LRI$ and $\MQG$ is essential in dynamic problems---is an open
question we intend to settle. \Claude{This last sentence states the
research question implied by your outline; if you already have a
counterexample showing MQG is \emph{not} sufficient in the dynamic
model, it belongs here and would sharpen the paper's contribution
considerably. Please advise.}

Two special cases deserve mention. First, when there are only two
states, every pair of MLR signal processes ordered by any of the above
criteria is comparable in a particularly transparent way;
\citet{jewitt2007information} shows that, under the MLR property,
Lehmann's order is equivalent to Blackwell dominance holding on all
dichotomies (all two-state restrictions) of the model. Second, your
outline records the claim that in models with \emph{two actions}, the
four orders coincide under the MLR property.
\Claude{I have not found this two-action equivalence stated in Kim
(2023)---his results connect MQG, Lehmann, and Blackwell but not in a
form specific to binary-action problems; the dichotomy result of
Jewitt (2007) concerns two \emph{states}, not two actions. If the
two-action claim is a result you intend to prove, I will set it up as
a proposition with a proof to be supplied; please clarify the intended
statement (e.g., ``for two-action monotone decision problems, $\iX$
yields a higher value than $\iY$ for all priors and payoffs if and
only if $\iX \MQG \iY$'') and its provenance.}
